\documentclass[11pt,letterpaper]{article}
\usepackage{graphicx}
\usepackage{geometry}
\usepackage{amsmath}
\usepackage{amssymb}
\usepackage{hyperref}
\usepackage{url}
\usepackage{natbib}
\usepackage{booktabs}
\usepackage{xcolor}
\usepackage{microtype}
\usepackage{array}
\usepackage{multirow}

\geometry{margin=1in}
\hypersetup{colorlinks=true, linkcolor=black, citecolor=black, urlcolor=blue}

\title{\textbf{Abliteration Cannot Erase RLHF: Geometric Evidence from\\
SAE Encoder Directions in Qwen3-8B}}

\author{AMGrobelnik}

\date{}

\begin{document}

\maketitle

\begin{abstract}
Refusal ablation (abliteration) edits a safety-aligned model's weights to eliminate compliance
refusals by removing a single residual-stream direction. The standard narrative frames abliteration
as undoing RLHF safety fine-tuning. We test this at the encoder-direction level for Qwen3-8B using
the frozen Qwen-Scope base SAE (65,536 TopK features), comparing three checkpoints---base,
RLHF-instruct, and a community-released abliterated variant---across a balanced corpus of 500 harmful
and 500 benign prompts. Because SAE reconstruction quality at layer 22 falls below acceptable
thresholds ($R^2 \approx 0.51$ for instruct and abliterated variants), all metrics operate on
pre-topk encoder projections ($\mathbf{z}_{\text{pre}}$)---dense continuous projections onto all
encoder directions---rather than the post-topk sparse features validated as semantically interpretable
in single-model settings. Within this geometric regime, we introduce two per-feature metrics: the
Reversal Score $R(f)$, measuring directional preservation of RLHF-induced encoder-direction changes
under abliteration, and the OLS Slope $\beta_f$, measuring magnitude preservation. For Qwen3-8B,
and conditioned on the abliterated checkpoint faithfully implementing the refusal-direction removal,
every one of the top-500 RLHF-sensitive encoder directions has a positive Reversal Score
(mean $\bar{R} = 0.921 \pm 0.040$, $\min = 0.684$; zero directions reversed), and the OLS slope
is near unity ($\bar{\beta} = 0.986 \pm 0.176$; 86.4\% of directions with $\beta_f > 0.8$). A
noise-free mid-SC baseline (Cohen's $d = 1.087$) confirms selective preservation. Direct measurement
of refusal-direction geometry---mean $|\cos\theta| = 0.019$ between the refusal direction and
top-500 RLHF-modified encoder directions at layer 22---provides the mechanistic explanation:
abliteration's orthogonalization is geometrically incapable of touching these directions. A 64.6\%
overlap with benign-prompt RLHF-sensitive directions reveals that the preserved modifications span
both harm-specific and general RLHF stylistic changes. These findings establish, for Qwen3-8B, that
abliteration removes the compliance mechanism while leaving the representational modifications of
RLHF intact at approximately unit magnitude, pending validation of the pre-topk proxy and independent
verification of abliteration depth.
\end{abstract}

\section{Introduction}

Refusal ablation---commonly called \emph{abliteration}---edits a safety-aligned language model's
weights to eliminate its compliance refusals \citep{Arditi2024}. The technique identifies a
\emph{refusal direction} in the residual stream: the mean-difference vector between activations on
harmful versus harmless prompts. All weight matrices that write to the residual stream are then
orthogonalized against this direction, producing a model that complies with previously-refused
requests without retraining. The standard framing treats abliteration as \emph{undoing} RLHF safety
alignment: if RLHF installs a refusal mechanism, abliteration removes it, restoring a base-like model.

This framing has direct implications for AI safety. If abliteration truly erases RLHF modifications,
the safety imprint is fragile---a single weight-editing step with minimal compute removes it entirely.
If abliteration only removes the \emph{behavioral expression} of safety training while leaving
\emph{representational} modifications intact, the picture differs substantially: the model's internal
representations may still carry RLHF-induced structure, providing a substrate for lightweight safety
restoration without full retraining. Distinguishing these scenarios matters both for threat modeling
and for designing targeted safety-recovery interventions.

Two geometric studies establish that harmful-prompt representations survive abliteration: LatentBiopsy
\citep{LlorenteSaguer2026a} finds that a linear classifier on Qwen model residual streams retains
AUROC $\geq 0.937$ after abliteration, and a cross-family companion study \citep{LlorenteSaguer2026b}
replicates this across 12 variants (Qwen, Llama, Gemma; 0.5B--9B parameters) with mean AUROC 0.982
surviving within $\pm 0.003$. Something in the internal representation persists---but holistic AUROC
collapses all representational information into a single detection score and cannot answer \emph{which
specific encoder-direction structures survive, at what magnitude, or why}. Answering these mechanistic
questions requires a feature-level vocabulary.

Sparse Autoencoders (SAEs) \citep{Elhage2022,Bricken2023,Gao2024,Templeton2024} decompose
residual-stream activations into tens of thousands of near-monosemantic features, providing exactly
this vocabulary. Prior SAE-based analyses of refusal identify which features \emph{drive} refusal
behavior within a single model \citep{Yeo2025,Kim2026}. Our work addresses a complementary question:
which encoder directions are modified by RLHF, and do those modifications survive abliteration in
direction \emph{and} magnitude? A key methodological complication is that reconstruction quality of
a base-trained SAE degrades on instruct-variant activations \citep{Chopra2026,Li2025}, ruling out
the standard post-topk sparse feature representation at the layers with the richest RLHF signal. We
address this by working with \emph{pre-topk encoder projections} ($\mathbf{z}_{\text{pre}}$)---dense
continuous projections of residual-stream activations onto all 65,536 SAE encoder directions---and
introduce two metrics interpretable in this geometric regime. The pre-topk proxy is validated
indirectly via the orthogonality result (Section~\ref{sec:orthogonality}), which is independent of
the sparse/dense distinction, and a direct $\mathbf{z}_{\text{pre}}$/$\mathbf{z}_{\text{post}}$
comparison is deferred to future work.

We apply the frozen Qwen-Scope base SAE (\texttt{Qwen/SAE-Res-Qwen3-8B-Base-W64K-L0\_50})
\citep{Deng2026} across three Qwen3-8B checkpoints---base, RLHF-instruct, and a community-released
abliterated variant---on a balanced corpus of 500 harmful and 500 benign prompts from seven datasets.
The \emph{Reversal Score} $R(f)$ measures whether abliteration preserves or reverses the direction of
the RLHF-induced change to encoder direction $f$. The \emph{OLS Slope} $\beta_f$ measures whether
abliteration preserves the magnitude of that change. For Qwen3-8B, conditioned on abliteration
completeness, every top-500 RLHF-sensitive encoder direction has a positive Reversal Score
(mean $\bar{R} = 0.921$, $\min = 0.684$, zero directions reversed), and the OLS slope is near unity
($\bar{\beta} = 0.986 \pm 0.176$, 86.4\% of directions with $\beta_f > 0.8$). Direct orthogonality
measurement (mean $|\cos\theta| = 0.019$ at layer 22) provides the geometric mechanism. A 64.6\%
overlap with benign-prompt RLHF-sensitive directions establishes that the preserved modifications are
not exclusively harm-specific.

\medskip
\noindent\textbf{Summary of Contributions:}
\begin{itemize}
  \item We introduce the \emph{Reversal Score}, a per-encoder-direction Pearson correlation confirming
    directional preservation of RLHF-modified SAE encoder directions under abliteration for Qwen3-8B
    ($\bar{R} = 0.921$ at layer 22, zero reversed directions; Section~\ref{sec:reversal}).%
    \footnote{Code: \url{https://github.com/AMGrobelnik/ai-invention-4ac301-abliteration-cannot-erase-rlhf-geometric/tree/main/round-2/experiment-1}}

  \item We introduce the \emph{OLS Slope} $\beta_f$, a companion magnitude metric, confirming that
    abliteration preserves not just direction but approximately full amplitude of RLHF-induced
    encoder-direction changes ($\bar{\beta} = 0.986 \pm 0.176$ at layer 22; Section~\ref{sec:ols}).%
    \footnote{Code: \url{https://github.com/AMGrobelnik/ai-invention-4ac301-abliteration-cannot-erase-rlhf-geometric/tree/main/round-4/experiment-1}}

  \item We validate selective preservation against a noise-free mid-SC baseline
    ($\text{SC}_{\text{median}} = 1.12 \gg 0.01$): top-500 $\bar{R} = 0.921$ vs.\ mid-500
    $\bar{R} = 0.871$, Cohen's $d = 1.087$, $p < 10^{-300}$ (Section~\ref{sec:baseline}).

  \item We directly measure refusal-direction orthogonality to RLHF-modified SAE encoder directions
    at layers 22 and 20 (mean $|\cos\theta| = 0.0185$ and $0.0209$, respectively), providing the
    geometric mechanism that explains zero reversed directions and $\bar{\beta} \approx 1$
    (Section~\ref{sec:orthogonality}).%
    \footnote{Code: \url{https://github.com/AMGrobelnik/ai-invention-4ac301-abliteration-cannot-erase-rlhf-geometric/tree/main/round-3/experiment-1}}

  \item We show that 64.6\% of the top-500 RLHF-sensitive encoder directions overlap with the
    top-500 benign-prompt RLHF-sensitive directions, establishing that the preserved RLHF
    modifications span general behavioral changes (output format, instruction style) in addition to
    harm-specific semantics (Section~\ref{sec:overlap}).
\end{itemize}

\section{Related Work}

\paragraph{Refusal ablation.}
\citet{Arditi2024} establish that refusal is mediated by a single residual-stream direction and that
orthogonalizing all weight matrices against it reliably eliminates compliance across 13 models up to
72B parameters. \citet{Joad2026} challenge the single-direction framing, demonstrating that refusal
behaviors correspond to geometrically distinct directions across eleven behavioral categories (safety
refusals, excessive refusal, incomplete-request refusals, etc.), though they find that linear steering
along any refusal-related direction produces nearly identical refusal-to-over-refusal trade-offs. This
finding is directly relevant to our orthogonality measurement (Section~\ref{sec:orthogonality}): the
mean $|\cos\theta| = 0.019$ between the mean-difference refusal vector and RLHF-modified encoder
directions is a necessary but not sufficient condition for explaining zero reversed features if
abliteration removes a multi-dimensional subspace rather than a single direction. We address this in
Section~\ref{sec:orthogonality} and treat the multi-dimensional extension as explicit future work.
\citet{Cristofano2026} show the raw refusal vector is polysemantic, entangling safety circuitry with
style; spectral residualization reduces collateral KL divergence while preserving abliteration.
\citet{Nanfack2026} propose an optimal-transport alternative reducing distributional shift. Our work
examines what all these variants share: the refusal direction (or subspace), however constructed,
leaves RLHF-modified SAE encoder directions intact in both direction and magnitude.

\paragraph{SAEs applied to refusal analysis.}
\citet{Yeo2025} use SAEs trained on pretraining data to investigate refusal mechanisms in chat-tuned
models, finding that base-trained SAEs fail to cleanly encode the refusal circuitry of chat
models---supporting our observation that base SAE reconstruction quality degrades on instruct
activations ($R^2 \approx 0.51$ at layer 22). Their finding contextualizes our pre-topk proxy choice:
the encoder directions of a base SAE are not guaranteed to align with the semantically active features
of an instruct model, which is precisely why we restrict claims to \emph{encoder-direction geometry}
rather than asserting preservation of semantically validated sparse features. \citet{Kim2026} introduce
CRaFT, which uses cross-layer transcoders to map model computations into sparse feature circuits,
quantifying inter-feature influences on refusal-output logits rather than treating features as isolated
signals. CRaFT and our approach are complementary: CRaFT identifies which features causally govern
refusal outputs in a single model; our Reversal Score and OLS Slope identify which encoder directions
RLHF modified and whether abliteration erases those modifications. Both are feature-level analyses,
but they use different feature vocabularies (transcoder features vs.\ frozen base SAE encoder
directions) and answer different questions (causal control of refusal vs.\ survival of RLHF
modifications under abliteration).

\paragraph{Geometry of harmful intent surviving abliteration.}
LatentBiopsy \citep{LlorenteSaguer2026a} and the companion geometry study \citep{LlorenteSaguer2026b}
are the direct predecessors establishing the phenomenon this work explains mechanistically. Their key
limitation---collapsing representational information into a single AUROC score---is precisely the gap
our encoder-direction analysis addresses. The OLS slope results ($\bar{\beta} = 0.986$) and direct
orthogonality measurement ($|\cos\theta| = 0.0185$) provide the mechanistic account that a holistic
detection score cannot.

\paragraph{SAEs as cross-variant measurement instruments.}
\citet{Chopra2026} uses frozen base SAEs as cross-variant diagnostics for supervised fine-tuning
(SFT), finding that SAE latent cosine similarity drops to 0.509--0.557 at layer 22 after SFT while
raw activation cosine stays above 0.96---validating the frozen base-SAE approach and motivating the
pre-topk encoder fallback when reconstruction quality degrades. \citet{Jiralerspong2026} introduce
crosscoders---SAEs trained jointly on two or more model variants---that partition features into shared
and model-specific categories without requiring a pre-existing vocabulary. \citet{Kassem2026} extend
this to narrow fine-tuning regimes with Delta-Crosscoder. Our frozen base-SAE approach is
complementary: it requires zero additional training, reuses Qwen-Scope's existing feature vocabulary,
and extends naturally to the three-model base$\to$instruct$\to$abliterated trajectory. \citet{Li2025}
demonstrate an $\approx$8$\times$ MSE gap between base-SAE reconstruction on instruct versus
purpose-trained instruct SAE activations, directly motivating our reconstruction quality gate and
pre-topk fallback. \citet{Shportko2026} localize RL-induced tool-calling in Qwen2.5-3B to a compact
steerable crosscoder feature set---the complement of our finding about which encoder directions
abliteration \emph{cannot} remove.

\paragraph{Safety fine-tuning mechanisms.}
\citet{Jain2024} establish that safety fine-tuning projects unsafe activations into MLP null spaces,
explaining adversarial circumvention. Our finding extends this picture: at the SAE encoder level,
RLHF-modified encoder directions are near-orthogonal to the refusal direction, so abliteration---which
removes the refusal direction---cannot erase them geometrically.

\section{Methodology}

\subsection{Three-Model Setup}
\label{sec:setup}

All experiments use three Qwen3-8B checkpoints sharing identical 36-layer architecture
($d_{\text{model}} = 4096$, bfloat16, transformers $\geq$ 4.51.0).%
\footnote{Code: \url{https://github.com/AMGrobelnik/ai-invention-4ac301-abliteration-cannot-erase-rlhf-geometric/tree/main/round-1/research-1}}

\noindent\emph{Base}: \texttt{Qwen/Qwen3-8B-Base}---the unaligned pretrained model, providing the
representational baseline against which RLHF-induced changes are measured.

\noindent\emph{Instruct}: \texttt{Qwen/Qwen3-8B} with \texttt{enable\_thinking=False}---the
RLHF-aligned model. Chain-of-thought suppression isolates safety-relevant activations on the final
input token.

\noindent\emph{Abliterated}: \texttt{huihui-ai/Huihui-Qwen3-8B-abliterated-v2} (v2 corrects a
layer-0 output-garbling artifact present in v1). The exact layer set and prompt set used for the
refusal direction computation are not publicly documented. The findings in this paper are conditioned
on this checkpoint faithfully implementing a refusal-direction removal: if abliteration is incomplete,
the high Reversal Scores and OLS slopes would trivially reflect the abliterated model being close to
the instruct model by construction rather than preservation under genuine erasure.
Section~\ref{sec:orthogonality} and the Limitations section discuss this constraint.

All models are accessed via \texttt{model.model.layers} for residual-stream hook injection at the
final token position of each input.

\subsection{SAE Measurement Instrument}
\label{sec:sae}

The frozen Qwen-Scope SAE \texttt{Qwen/SAE-Res-Qwen3-8B-Base-W64K-L0\_50} serves as the shared
measurement instrument \citep{Deng2026}. A critical architectural note: Qwen-Scope implements
\emph{TopK} architecture ($k = 50$ features retained via \texttt{torch.topk}), not JumpReLU. For a
residual activation $\mathbf{x} \in \mathbb{R}^{4096}$:

\begin{align}
  \mathbf{z}_{\text{pre}} &= \mathbf{x}W_{\text{enc}}^\top + \mathbf{b}_{\text{enc}}
    \;\in\; \mathbb{R}^{65536} \\
  \mathbf{z} &= \text{TopK}(\mathbf{z}_{\text{pre}},\; k{=}50)
\end{align}

\noindent\textbf{Pre-topk proxy.} The reconstruction quality gate (Section~\ref{sec:recon}) fails
at layer 22: $R^2 = 0.574$ (base), $0.512$ (instruct), $0.512$ (abliterated)---all below the 0.85
threshold \citep{Li2025}. All feature metrics therefore use $\mathbf{z}_{\text{pre}}$ rather than the
post-topk sparse representation. The pre-topk vector is dense by construction: every prompt receives
a non-zero value on all 65,536 encoder directions simultaneously, so no single direction is ``off''
in the interpretability sense. The metrics we compute (Reversal Score, OLS Slope) measure
\emph{geometric properties of encoder directions in a continuous projection space}, not activation of
near-monosemantic features in the interpretability sense. The orthogonality result
(Section~\ref{sec:orthogonality}) is fully independent of the sparse/dense distinction---it is a
direct measurement of the angle between weight vectors---and remains the most secure finding. A direct
comparison of top-500 feature rankings and Reversal Score rank correlations between $\mathbf{z}_{\text{pre}}$
and $\mathbf{z}_{\text{post}}$ on the highest-$R^2$ subset is an important validation deferred to
future work. If such a comparison shows Spearman $\rho > 0.7$ between $\mathbf{z}_{\text{pre}}$ and
$\mathbf{z}_{\text{post}}$ Reversal Scores, the pre-topk proxy is defensible as a surrogate; if it
diverges, the encoder-direction findings must be interpreted as purely geometric properties without
semantic grounding.

\subsection{Feature Metrics}
\label{sec:metrics}

For each SAE encoder direction $f$ and $N = 500$ harmful prompts $\{p_i\}_{i=1}^N$, let
$a_m(f, p_i)$ denote the pre-topk encoder activation of direction $f$ for model
$m \in \{\text{base}, \text{inst}, \text{abl}\}$.

\noindent\textbf{Safety Change} quantifies the magnitude of the RLHF modification to encoder
direction $f$:
\begin{equation}
  \text{SC}(f) = \frac{1}{N}\sum_{i=1}^N |a_{\text{inst}}(f, p_i) - a_{\text{base}}(f, p_i)|
\end{equation}

Encoder directions are ranked by $\text{SC}(f)$: the top-500 form the \emph{RLHF-sensitive set}.
For null-baseline analysis, a \emph{mid-500} group is drawn uniformly from the 40th--60th SC
percentile.

\noindent\textbf{Reversal Score} quantifies whether abliteration reverses or preserves the
\emph{direction} of the RLHF modification to encoder direction $f$:
\begin{align}
  \boldsymbol{\Delta}_{\text{inst}}(f) &= [a_{\text{inst}}(f, p_i) - a_{\text{base}}(f, p_i)]_{i=1}^N \\
  \boldsymbol{\Delta}_{\text{abl}}(f)  &= [a_{\text{abl}}(f, p_i)  - a_{\text{base}}(f, p_i)]_{i=1}^N \\
  R(f) &= \text{PearsonR}(\boldsymbol{\Delta}_{\text{inst}}(f),\; \boldsymbol{\Delta}_{\text{abl}}(f))
\end{align}

$R(f) \approx +1$ indicates abliteration preserves the RLHF modification (\emph{imprinted});
$R(f) \approx -1$ indicates abliteration reverses it (\emph{reversed}). This is a correlation over
prompt-level scalar activations, not a cosine similarity of scalars---ensuring that structure in the
result is an empirical finding, not an artifact of the definition.

\noindent\textbf{OLS Slope} resolves the scale-invariance of the Reversal Score. A feature where
abliteration retains only 5\% of the RLHF-induced shift in magnitude but preserves its direction
would still yield $R(f) \approx 1.0$. The no-intercept ordinary least-squares slope:
\begin{equation}
  \beta_f = \frac{\boldsymbol{\Delta}_{\text{inst}}(f)^\top \boldsymbol{\Delta}_{\text{abl}}(f)}
                 {\boldsymbol{\Delta}_{\text{inst}}(f)^\top \boldsymbol{\Delta}_{\text{inst}}(f)}
\end{equation}

$\beta_f \approx 1$ means the abliterated model's encoder-direction activations shift by the same
magnitude as the instruct model's, relative to base. $\beta_f \approx 0$ means only the directional
pattern survives with negligible magnitude. $\beta_f > 1$ indicates overshooting.

\subsection{Benign-SC Overlap}
\label{sec:benign}

To assess whether the top-500 RLHF-sensitive encoder directions reflect \emph{harm-specific}
modifications or \emph{general RLHF} modifications (style, tone, response format), we separately
compute $\text{SC}_{\text{benign}}(f)$ using the 500 benign prompts and report the overlap between
the top-500 by harmful-SC and the top-500 by benign-SC.

\subsection{Layer Selection and Reconstruction Quality Gate}
\label{sec:recon}

\textbf{Layer selection.} We compute $\text{SC}(f)$ variance across all 65,536 encoder directions at
candidate layers \{14, 20, 22\} on 50 harmful prompts. Layer 22 is selected (variance 0.1721 vs.\
0.0547 at layer 20 and 0.0127 at layer 14); all primary results use layer 22 with layer 20 as
replication.

\textbf{Reconstruction quality gate.} At layer 22, $R^2 = 0.574$ (base), $0.512$ (instruct),
$0.512$ (abliterated)---all below the 0.85 threshold \citep{Li2025}---triggering the pre-topk
encoder-only analysis path.

\subsection{Semantic Labeling Protocol}
\label{sec:labeling}

For the top-100 encoder directions by $\text{SC}(f)$, we identify the 5 prompts with the highest
pre-topk encoder activation and provide the first 300 characters of each to Llama-3.1-8B-Instruct
via OpenRouter. The cross-family labeler avoids within-Qwen circularity. Each direction is classified
as \texttt{harm\_concept}, \texttt{safety\_refusal}, \texttt{general\_capability}, or \texttt{other}.
Results are treated as qualitative observations only: the labeling uses character-prefix truncation of
most-activating prompts rather than per-token activation windows; labels reflect prompt topic more
than token-level representations; and no Cohen's $\kappa$ validation against a manual gold set was
performed.

\section{Evaluation Corpus}
\label{sec:corpus}

The evaluation corpus comprises 500 harmful and 500 benign prompts from seven publicly licensed
datasets.

\noindent\textbf{Harmful sources:} BeaverTails \citep{Ji2023} (176 prompts; NeurIPS 2023,
human-labeled QA spanning 14 harm categories), AdvBench (98 prompts; standard adversarial behavior
dataset used in abliteration research \citep{Arditi2024}), TrustAIRLab in-the-wild jailbreak prompts
(86 prompts; real-world jailbreaks from ACM CCS 2024), HarmBench \citep{Mazeika2024} (85 prompts;
covering chemical, cybercrime, physical harm, and illegal categories), and
JailbreakBench/JBB-Behaviors \citep{Chao2024} (55 prompts).

\noindent\textbf{Benign sources:} Jailbreak-classification (151 prompts; labeled benign user
requests), JBB-Behaviors benign counterparts \citep{Chao2024} (46 prompts), Toxic-Chat (104 prompts;
real user prompts labeled non-toxic), TrustAIRLab regular prompts (93 prompts), and Dolly-15k
\citep{Conover2023} (106 prompts; human-written instruction-following tasks).

All sources are processed through MinHash LSH deduplication (threshold 0.5, 128 permutations,
character 5-gram shingling), token-length filtering ($\leq$512 tokens), and category capping
($\leq$35\% per harm subcategory). The final corpus contains 1,000 prompts (500 harmful, 500 benign).

\section{Results}

\subsection{Layer Selection}
\label{sec:layers}

Layer 22 exhibits the highest Safety Change variance (0.1721), confirming it as the primary analysis
layer. The $R^2$ reconstruction gate at layer 22 yields values below 0.85 for all three model
variants (base: 0.574, instruct: 0.512, abliterated: 0.512), triggering the pre-topk encoder-only
analysis path (Section~\ref{sec:sae}).

\subsection{Reversal Score Distribution}
\label{sec:reversal}

Every top-500 RLHF-sensitive encoder direction has a positive Reversal Score
(mean $= 0.921$, std $= 0.040$, min $= 0.684$, max $= 0.989$). No direction has a negative or
near-zero score. Hartigan's dip test yields statistic $= 0.0123$, $p = 0.930$---strongly unimodal.
Spearman correlation mean $= 0.913$, confirming rank-based robustness.

Table~\ref{tab:reversal} summarizes the count of reversed and imprinted directions at all tested
thresholds. Zero directions are reversed at any threshold.

\begin{table}[!htbp]
  \centering
  \caption{Imprinted and reversed encoder-direction counts at all tested thresholds.}
  \label{tab:reversal}
  \begin{tabular}{lcc}
    \toprule
    Threshold & Imprinted ($R >$ threshold) & Reversed ($R < -$threshold) \\
    \midrule
    0.30 & 500 / 500 (100.0\%) & 0 / 500 (0\%) \\
    0.50 & 500 / 500 (100.0\%) & 0 / 500 (0\%) \\
    0.70 & 499 / 500 (99.8\%)  & 0 / 500 (0\%) \\
    0.95 & 125 / 500 (25.0\%)  & 0 / 500 (0\%) \\
    \bottomrule
  \end{tabular}
\end{table}

\subsection{OLS Slope: Magnitude Preservation}
\label{sec:ols}

A key limitation of the Reversal Score is scale-invariance: $R(f) \approx 1$ is consistent with
abliteration preserving only 5\% of the RLHF-induced magnitude while retaining the direction pattern.
The OLS slope $\beta_f$ resolves this ambiguity.

The $\beta_f$ distribution for the top-500 RLHF-sensitive encoder directions at layer 22 is tightly
concentrated near 1: mean $\bar{\beta} = 0.986 \pm 0.176$, median $= 0.985$, range $[-0.055,\, 1.725]$.
Specifically, 86.4\% of directions have $\beta_f > 0.8$ (abliteration preserves more than 80\% of
the RLHF-induced shift in magnitude), 45.4\% have $\beta_f > 1.0$ (slight amplification), and only
0.2\% are negative. \textbf{MAGNITUDE PRESERVED.} Within the pre-topk geometric regime, abliteration
does not attenuate the RLHF imprint---it preserves approximately unit magnitude.

Layer-20 replication: $\bar{\beta} = 0.954 \pm 0.191$, 80.2\% of directions with $\beta_f > 0.8$,
zero negative values---consistent with the layer-22 primary result.

\subsection{Null-Baseline Validation: Noise-Free Mid-SC Comparison}
\label{sec:baseline}

A potential concern with a bottom-500 null baseline is that near-zero Safety Change values may make
PearsonR between near-zero vectors potentially unreliable. We address this with a noise-free mid-SC
baseline: 500 directions uniformly sampled from the 40th--60th SC percentile, with
$\text{SC}_{\text{median}} = 1.12$ (versus the top-500 minimum of 2.79).

At layer 22: top-500 $\bar{R} = 0.921$, mid-500 $\bar{R} = 0.871$ (Cohen's $d = 1.087$,
two-sample $t$-test $t = 17.2$, $p < 10^{-300}$). \textbf{MID-SC CONFIRMS SELECTIVE PRESERVATION.}
The top-500 encoder directions show significantly higher Reversal Scores than mid-SC controls that
have non-trivial RLHF sensitivity and noise-free $\boldsymbol{\Delta}_{\text{inst}}$ vectors,
confirming that the high Reversal Scores reflect genuine selective preservation rather than any
measurement artifact.

\subsection{Direct Measurement of Refusal-Direction Orthogonality}
\label{sec:orthogonality}

The mechanistic interpretation of zero reversed directions and $\bar{\beta} \approx 1$ requires a
testable prediction: the refusal direction $\mathbf{r}$ must be nearly orthogonal to all
RLHF-modified SAE encoder directions. We test this directly.

The refusal direction $\mathbf{r}$ is computed as the mean difference of instruct residual-stream
activations on harmful versus harmless prompts at layer 22. For each of the top-500 RLHF-sensitive
encoder directions $f$:
\begin{equation}
  \cos\theta_f = \frac{\mathbf{r} \cdot \mathbf{w}_{\text{enc}}(f)}
                      {\|\mathbf{r}\| \|\mathbf{w}_{\text{enc}}(f)\|}
\end{equation}

\noindent\textbf{Layer 22:} mean $|\cos\theta_f| = 0.01854$ (authoritative value; 0.019 in the
abstract is rounded). \textbf{ORTHOGONALITY CONFIRMED.} A cosine of 0.0185 means the refusal
direction projects onto less than 1.9\% of each RLHF-modified encoder direction's magnitude.
\textbf{Layer 20:} mean $|\cos\theta_f| = 0.02086$. Orthogonality confirmed at both layers;
inter-layer agreement (matching to 5 decimal places across independent experimental runs) confirms
these values are stable.

This orthogonality result is the most secure finding in this paper: it is a direct measurement of the
angle between weight vectors and is fully independent of the pre-topk/post-topk distinction, the
abliteration checkpoint's completeness, or any activation-level measurement.

This orthogonality geometrically explains the $\bar{\beta} \approx 1$ result: abliteration
orthogonalizes all weight matrices against $\mathbf{r}$. Because $\mathbf{r}$ is nearly orthogonal
to all RLHF-modified SAE encoder directions, the orthogonalization cannot substantially change how
any RLHF-modified encoder direction's activations are computed---in direction or in magnitude.

\noindent\textbf{Caveat on multi-dimensional refusal.} \citet{Joad2026} demonstrate that refusal
spans geometrically distinct directions across eleven behavioral categories. If abliteration removes
a multi-dimensional refusal \emph{subspace} rather than a single vector, the orthogonality argument
requires that the full ablation subspace be near-orthogonal to all RLHF-modified encoder directions.
The current measurement tests only the mean-difference vector $\mathbf{r}$. An extension to the top-3
principal components of the (harmful $-$ harmless) activation differences would determine whether
this argument holds for higher PCs; this is deferred to future work. If a higher PC shows mean
$|\cos\theta_f| \geq 0.05$ against the top-500 encoder directions, the mechanistic explanation would
require qualification.

\subsection{Benign-SC Overlap: What the Preserved Imprint Represents}
\label{sec:overlap}

We compute $\text{SC}_{\text{benign}}(f)$---Safety Change evaluated on 500 benign prompts---for all
65,536 SAE encoder directions, then measure the overlap between the top-500 by harmful-SC and the
top-500 by benign-SC.

\textbf{Overlap count: 323 / 500 = 64.6\%. STYLE-DOMINANT verdict.} A majority of the top-500
RLHF-sensitive encoder directions also appear in the benign-SC top-500, indicating that RLHF modified
these directions similarly when processing harmful and benign prompts. Approximately 64.6\% of the
preserved imprint captures general RLHF modifications (response formality, output structure,
instruction-following style) rather than harm-specific semantic content. The remaining 35.4\% (176
directions) are uniquely sensitive to harmful prompts and more plausible carriers of harm-specific
semantic representations.

This finding contextualizes the OLS slope result ($\bar{\beta} = 0.986$): abliteration
comprehensively preserves all RLHF encoder-direction modifications regardless of whether they encode
harm-specific or general-style content. The near-orthogonality of the refusal direction to all
RLHF-modified encoder directions means abliteration cannot selectively erase harm-semantic changes
while leaving style changes intact---it fails to erase either category because all reside in a
subspace near-orthogonal to $\mathbf{r}$.

\subsection{Centroid Analysis with Control Subspaces}

To validate that the preserved imprint is specific to the RLHF-sensitive subspace, we compute
centroid distance ratios $d_{\text{abl-base}} / d_{\text{abl-inst}}$ (cosine) for three
encoder-direction subspaces.

\noindent\textbf{Layer 22 (harmful prompts):} Subspace A (top-500 RLHF-sensitive) ratio $= 4.35\times$,
Subspace B (bottom-500 RLHF-insensitive) ratio $= 1.16\times$, Subspace C (random-500) ratio
$= 2.37\times$. \textbf{Layer 20:} Subspace A ratio $= 7.31\times$, Subspace B ratio $= 1.01\times$,
Subspace C ratio $= 1.63\times$ (verdict: CENTROID SELECTIVE). The RLHF-sensitive subspace shows
4--7$\times$ greater instruct-closeness than the RLHF-insensitive control (B $\approx 1.0$--$1.2$
vs.\ A $= 4.35$--$7.31$), confirming that the imprint is specific to the directions RLHF modified,
not a global consequence of model similarity. Layer-20 Subspace B reaching $1.01\times$ (essentially
no instruct-closeness in the RLHF-insensitive subspace) provides the cleanest control.

\subsection{Downstream Classification Accuracy}

This null result follows directly from the 64.6\% benign-SC overlap (Section~\ref{sec:overlap}): the
top-500 RLHF-sensitive encoder directions are elevated by RLHF on \emph{both} harmful and benign
prompts, eliminating prompt-level discriminability by construction. This distinguishes \emph{what
RLHF modified} (broad encoder-direction changes) from \emph{what selectively activates on harmful
versus benign prompts at inference} (a different, narrower set of encoder directions not analyzed in
this work).

A per-prompt harm score (mean pre-topk activation across all top-500 RLHF-sensitive encoder
directions with median threshold) achieves \textbf{0.476 accuracy}---near-chance on the balanced
corpus---versus a keyword baseline of \textbf{0.630}. This null classification result confirms the
64.6\% benign-SC overlap from an independent angle: RLHF-sensitive encoder directions elevated on
harmful prompts are similarly elevated on benign prompts, yielding no discriminative signal.

\subsection{Semantic Labeling}

Llama-3.1-8B-Instruct assigned labels to 100 top-SC encoder directions: the majority received
\texttt{harm\_concept} labels, with smaller proportions of \texttt{safety\_refusal},
\texttt{general\_capability}, and \texttt{other}. These results are qualitative observations only.
The labeling uses character-prefix truncation of most-activating prompts rather than per-token
activation windows; no Cohen's $\kappa$ validation against a manual gold set was performed; and the
64.6\% benign-SC overlap implies that many stylistic directions receive \texttt{harm\_concept} labels
because they activate most strongly on harmful-prompt contexts rather than encoding harm-specific
semantic content. Quantitative reliance on the specific labeling percentages requires a manual
gold-set annotation with $\kappa \geq 0.6$, deferred to future work.

\subsection{Layer-20 Replication}

Table~\ref{tab:replication} summarizes full quantitative replication at layer 20.

\begin{table}[!htbp]
  \centering
  \caption{Quantitative results at layers 22 and 20.}
  \label{tab:replication}
  \begin{tabular}{lcc}
    \toprule
    Metric & Layer 22 & Layer 20 \\
    \midrule
    Top-500 Reversal Score ($\bar{R}$)   & 0.921 & 0.924 \\
    Mid-500 Reversal Score ($\bar{R}$)   & 0.871 & 0.858 \\
    Cohen's $d$ (top vs.\ mid)           & 1.087 & --- \\
    Top-500 OLS Slope ($\bar{\beta}$)    & $0.986 \pm 0.176$ & $0.954 \pm 0.191$ \\
    Fraction $\beta_f > 0.8$             & 86.4\% & 80.2\% \\
    Fraction $\beta_f < 0$               & 0.2\%  & 0.0\%  \\
    Mean $|\cos(\mathbf{r}, \mathbf{w}_{\text{enc}})|$ top-500 & 0.01854 & 0.02086 \\
    Centroid A (RLHF-sensitive)          & 4.35$\times$ & 7.31$\times$ \\
    Centroid B (RLHF-insensitive)        & 1.16$\times$ & 1.01$\times$ \\
    Centroid C (random-500)              & 2.37$\times$ & 1.63$\times$ \\
    Centroid verdict                     & MODERATE & SELECTIVE \\
    \bottomrule
  \end{tabular}
\end{table}

\noindent\textbf{REPLICATION CONFIRMED} at both layers.

\section{Discussion}

\subsection{Geometric Interpretation of Encoder-Direction Preservation}

The OLS slope result ($\bar{\beta} = 0.986$, 86.4\% of directions with $\beta_f > 0.8$) establishes
that the preservation is comprehensive in both direction and magnitude. The direct orthogonality
measurement ($|\cos\theta| = 0.0185$) provides the mechanistic account: abliteration removes
$\mathbf{r}$ from all weight matrices, but for abliteration to change the projection of a
residual-stream activation onto SAE encoder direction $f$ by even 10\%, $\mathbf{r}$ would need to
project onto $\mathbf{w}_{\text{enc}}(f)$ with $|\cos\theta| \geq 0.1$. The measured mean
$|\cos\theta| = 0.0185$ falls 5$\times$ below this threshold, explaining why both direction and
magnitude are preserved: the orthogonalization barely touches the RLHF-modified encoder directions.
This interpretation is consistent with \citeauthor{Cristofano2026}'s observation \citep{Cristofano2026}
that the raw refusal direction primarily captures stylistic compliance patterns.

We note that this mechanistic account applies in the pre-topk projection regime. Whether the same
geometry holds for the post-topk sparse activations validated as semantically interpretable in
single-model settings \citep{Bricken2023,Templeton2024} remains an open question. The orthogonality
result---which is a property of weight vectors, not activations---directly implies that \emph{any}
function of the residual stream that can be expressed as a projection onto $\mathbf{w}_{\text{enc}}(f)$
will be approximately unaffected by abliteration; the question is whether the semantically active
post-topk features are well-described by such projections, which requires a
$\mathbf{z}_{\text{pre}}$/$\mathbf{z}_{\text{post}}$ correlation analysis not yet available.

\subsection{What the Preserved Imprint Represents}

The 64.6\% benign-SC overlap requires careful interpretation. The majority of RLHF-sensitive encoder
directions are modified similarly when processing harmful and benign prompts, indicating these
directions capture \emph{general} RLHF behavioral changes (formality, helpfulness, instruction-following
structure) rather than harm-specific semantic content. The surviving 35.4\% uniquely-harmful
directions are the more plausible carriers of harm-concept semantic content.

Crucially, the near-orthogonality of the refusal direction to all RLHF-modified encoder directions
(harmful-specific and style-related alike) means abliteration geometrically cannot remove any category
of RLHF modification. The behavioral regression to base-like output arises because the compliance
mechanism (the refusal direction) is removed, not because the underlying encoder-direction
modifications are erased. The model forgets how to express refusal while retaining all
representational changes RLHF introduced---at approximately full magnitude.

This also explains the downstream classification null (accuracy $= 0.476$): the RLHF-sensitive
encoder directions are elevated on \emph{both} harmful and benign prompts, so mean-activation
discrimination is not available. The representational change RLHF introduced is a broad
encoder-direction shift, not a selective harm-detector.

\subsection{Implications for Safety Restoration}

If abliteration preserves the RLHF imprint at approximately unit magnitude, the challenge for
targeted safety restoration is not to amplify surviving encoder-direction modifications---they are
already at full instruct-model activation levels. The challenge is to restore the \emph{compliance
mechanism}: the geometric pathway from current activations to refusal outputs, implemented by the
refusal direction $\mathbf{r}$. Restoration requires re-injecting $\mathbf{r}$ into the weight
matrices without disturbing the already-intact encoder-direction modifications.

This reframes the restoration problem as geometric re-injection rather than semantic amplification.
The approach of \citet{Shportko2026} identifies compact steerable feature sets for RL-induced
capabilities; an analogous method targeting the encoder directions that mediate harm-concept-to-refusal
conversion could enable targeted re-injection. The 35.4\% uniquely-harmful encoder directions are the
most promising candidates for such targeted intervention, pending confirmation that they carry
semantically coherent harm-specific content (which requires $\mathbf{z}_{\text{post}}$ analysis).

\subsection{Limitations}

\textbf{Pre-topk encoder projections.} All feature metrics use $\mathbf{z}_{\text{pre}}$ (dense
encoder projections) rather than the post-topk sparse activations validated as interpretable by
\citet{Bricken2023} and \citet{Templeton2024}. Pre-topk values are dense by construction---every
encoder direction receives a non-zero value on every prompt---so findings are more accurately
characterized as \emph{geometric properties of encoder directions in a continuous projection space}
than as \emph{feature-level preservation of near-monosemantic SAE features}. A validation experiment
is required before semantic claims can be grounded: on the highest-$R^2$ prompt subset, compare
top-500 features selected by $\text{SC}(\mathbf{z}_{\text{pre}})$ versus
$\text{SC}(\mathbf{z}_{\text{post}})$ (overlap fraction) and compute Spearman $\rho$ between
$\mathbf{z}_{\text{pre}}$ and $\mathbf{z}_{\text{post}}$ Reversal Scores. If $\rho > 0.7$ and
overlap $> 60\%$, the pre-topk proxy is defensible; if they diverge, the semantic labeling results
must be interpreted purely as observations about encoder-direction geometry. This is the most critical
open validation gap.

\textbf{Abliteration depth verification.} The \texttt{huihui-ai/Huihui-Qwen3-8B-abliterated-v2}
checkpoint was not independently verified. If abliteration is incomplete---even partially---the high
Reversal Scores and $\bar{\beta} \approx 1$ could trivially reflect the abliterated model being close
to the instruct model by construction. Two measurements would resolve this: (1) fraction of AdvBench
or HarmBench prompts on which the abliterated checkpoint produces a non-refusal response (expected
near 100\% for complete abliteration); (2) mean
$|\text{Proj}_{\mathbf{r}}(\mathbf{h}_{\text{abl}}(p))| / |\mathbf{h}_{\text{abl}}(p)|$ at layer 22
for the 500 harmful prompts (expected near zero if $\mathbf{r}$ has been removed). These measurements
can be computed from already-cached activations. Until they are reported, the main findings should be
interpreted as conditional on abliteration completeness.

\textbf{Multi-dimensional refusal subspace.} The orthogonality measurement uses the mean-difference
vector $\mathbf{r}$ (a single vector). \citet{Joad2026} demonstrate that refusal spans geometrically
distinct directions across behavioral categories, implying abliteration may remove a multi-dimensional
subspace rather than a single vector. The mechanistic explanation of zero reversed directions requires
near-orthogonality of the \emph{full ablation subspace} to all RLHF-modified encoder directions, not
just the mean-difference summary. Computing mean $|\cos\theta_f|$ for the top-3 principal components
of the (harmful $-$ harmless) activation differences at layer 22 would determine whether the
orthogonality argument holds robustly; this is an important future measurement.

\textbf{Single model family and layers.} All results are for Qwen3-8B at layers 22 and 20. The
zero-reversed finding and $\bar{\beta} \approx 1$ result may not generalize to other model families
(Llama, Gemma) or other layer depths. The companion geometry studies \citep{LlorenteSaguer2026a,LlorenteSaguer2026b}
cover 12 variants, establishing the holistic survival phenomenon across families; analogous SAE-level
analysis for Llama-3.1-8B and Gemma-3-9B using GemmaScope and LlamaScope SAEs would test whether
encoder-direction preservation is a general property of the abliteration geometry.

\textbf{Semantic labeling without validation.} The qualitative observation that the majority of
top-SC encoder directions receive \texttt{harm\_concept} labels cannot support quantitative downstream
reasoning without Cohen's $\kappa$ validation against a manual gold set of at least 50 annotations.

\section{Conclusion}

We present the first encoder-direction-level mechanistic analysis of refusal abliteration using the
frozen Qwen-Scope base SAE across three Qwen3-8B variants. Two complementary per-direction
metrics---the Reversal Score (directional preservation) and the OLS Slope (magnitude preservation)---together
establish, for Qwen3-8B, a comprehensive geometric imprint: every top-500 RLHF-sensitive SAE encoder
direction has a positive Reversal Score (mean $\bar{R} = 0.921 \pm 0.040$, $\min = 0.684$) and an
OLS slope near unity ($\bar{\beta} = 0.986 \pm 0.176$, 86.4\% exceeding 0.8). A noise-free mid-SC
baseline (Cohen's $d = 1.087$) confirms selective preservation. Direct measurement of
refusal-direction geometry ($|\cos\theta| = 0.0185$ at layer 22, $0.0209$ at layer 20) provides the
mechanistic explanation: abliteration's orthogonalization is geometrically incapable of touching the
RLHF-modified encoder directions. A 64.6\% benign-SC overlap reveals that the preserved imprint
spans both harm-semantic and general RLHF stylistic encoder-direction modifications; abliteration
fails to erase any category because all reside in a subspace near-orthogonal to the refusal direction.
For Qwen3-8B, the abliterated model's behavior reverts to base-like, but its encoder-direction
projections remain instruct-like at full magnitude---a dissociation that reframes safety restoration
as geometric re-injection of the compliance direction rather than semantic amplification of surviving
features.

Three critical open gaps constrain full confidence in the broader interpretation: (1)
$\mathbf{z}_{\text{pre}}$/$\mathbf{z}_{\text{post}}$ validation (Spearman $\rho$ and feature-ranking
overlap at the highest-$R^2$ subset) to determine whether the geometric encoder-direction findings
correspond to semantically interpretable sparse features; (2) independent behavioral and geometric
verification of abliteration depth for the community checkpoint; and (3) extension of the
orthogonality measurement to the top-3 principal components of the refusal subspace, following
\citet{Joad2026}.

\noindent\textbf{Future directions:}
(i) $\mathbf{z}_{\text{pre}}$/$\mathbf{z}_{\text{post}}$ validation experiment;
(ii) abliteration depth verification via refusal-rate measurement and mean
$|\text{Proj}_{\mathbf{r}}|$ at layer 22;
(iii) multi-dimensional refusal subspace (top-3 PC orthogonality);
(iv) cross-family extension to Llama-3.1-8B and Gemma-3-9B using GemmaScope and LlamaScope SAEs;
(v) multi-layer analysis across all 36 Qwen-Scope layers;
(vi) investigation of the 35.4\% uniquely-harmful encoder directions as candidates for targeted
geometric restoration.

\bibliographystyle{plainnat}
\bibliography{references}

\end{document}
